\documentclass[12pt]{article}
\usepackage{amsfonts,amsthm,amsmath,amssymb,mathtools}
\usepackage{mathrsfs}
\usepackage{enumitem}
\usepackage{tikz-cd}
\usepackage{geometry}
\usepackage{mlmodern}
\usepackage{xcolor}

\geometry{letterpaper, portrait, margin=1in}

\setcounter{section}{0}

\renewcommand{\thesection}{\S\arabic{section}}
\renewcommand{\thesubsection}{\arabic{section}}
\newcommand{\wt}{\ensuremath{\operatorname{wt}}}
\newcommand{\LT}{\ensuremath{\text{\sc lt}}}
\newcommand{\Mat}{\ensuremath{\textsf{Mat}}}
\newcommand{\solution}{\noindent\textbf{Solution:~~}}
\newcommand{\SL}{\operatorname{SL}}
\newcommand{\GL}{\operatorname{GL}}
\newcommand{\imag}{\operatorname{im}}
\newcommand{\Orb}{\operatorname{Orb}}
\newcommand{\Stab}{\operatorname{Stab}}

\DeclareMathOperator{\lcm}{lcm}
\DeclareMathOperator{\im}{im}

\newtheorem{thm}{Theorem}
\newenvironment{theorem}[1]{
	\renewcommand{\thethm}{#1}
	\begin{thm}
		}{\end{thm}}

\newtheorem{lma}{Lemma}
\newenvironment{lemma}[1]{
	\renewcommand{\thelma}{#1}
	\begin{lma}
		}{\end{lma}}

\theoremstyle{definition}
\newtheorem{ex}{}
\newenvironment{exercise}[1]{
	\renewcommand{\theex}{#1}
	\begin{ex}
		}{\end{ex}}

\title{Homework 6}
\author{Connor Mooneyhan}
\date{October 27, 2025}

\begin{document}

\maketitle

Disclaimer: Generative AI was used in trying to understand the application of Sylow's theorems.

\begin{ex}
	Let $G$ be a finite group, let $p$ be a prime dividing $|G|$, let $H$ be a normal subgroup of $G$ and let $P$ be a Sylow
	$p$-subgroup of $H$.
	\begin{enumerate}
		\item Let $g \in G$.  Prove that there exists an $x \in H$ such that $gPg^{-1} = xPx^{-1}$.
		\item Use the previous part to conclude that $G = HN_G(P)$.
		\item Prove that $[G:H]$ divides $|N_G(P)|$.
	\end{enumerate}

\end{ex}

\begin{ex}
	Prove\footnote{Try rewriting a product of two transpositions as a product of $3$-cycles}
	that $A_n$ is generated by $3$-cycles.
\end{ex}
\begin{proof}
	Given two transpositions, they may ``overlap'' (i.e. have entries in common) in three different ways; they could be disjoint, they could share exactly one entry, or they could share both entries, in which case they are the same.

	If they are the same, then they are inverses of each other, and so their product is the identity, which is indeed generated by any 3-cycle since the order of a 3-cycle is 3.

	If they overlap in one entry only, then let the two transpositions be $(a\;b)$ and $(a\;c)$ without loss of generality, since we could always rearrange the matching entry to be first.
	Then $(a\;b)(a\;c)=(a\;c\;b)$, which \textit{is} a 3-cycle.

	Finally, if they are disjoint, then we can write the two transpositions as $(a\;b)$ and $(c\;d)$, and we see that \begin{equation*}
		(a\;c\;b)(a\;c\;d) = (a\;b)(c\;d),
	\end{equation*}
	so again it is generated by 3-cycles.

	Thus the identity is generated by 3-cycles, and every other even permutation is generated by 3-cycles since every even permutation is a product of pairs of transpositions, which we have shown are themselves generated by 3-cycles.
\end{proof}

\begin{ex}
	Prove that if $p$ is prime and $G$ is a nonabelian group of order $p^3$, then $Z(G) = [G,G]$.
	Show that $Z(G) \neq [G,G]$ for $G = D_{16}$.
\end{ex}
\begin{proof}
  Since $G$ is a $p$-group, $Z(G)$ is nontrivial, and thus $|Z(G)|$ is either $p$, $p^2$, or $p^3$.
  It cannot be $p^3$, since then $G$ would be abelian, and it cannot be $p^2$, because then $[G : Z(G)]$ would be $p$, making $G$ again abelian.
  So we must have $|Z(G)| = p$.
  Not sure where to go from there.
\end{proof}

\begin{ex}
	Let $G$ be a group of order $pqr$ where $p < q < r$ are primes.  Prove\footnote{Count elements!}
	that one of the Sylow subgroups of $G$ must be normal.

\end{ex}

\begin{ex}
	Prove that if $|G| \leq 100$ (or $< 168$ if you want more
	practice) and $G$ is simple, then either $|G| = 60$ or $|G| = p$ for
	$p$ a prime.  Many cases follow from problems on either this or the
	previous homework, and the remaining cases should follow either from
	counting elements via Sylow's theorem, or by using one of the
	permutation representations corresponding to a Sylow subgroup.  The
	cases $|G| = 90$ and $|G| = 144$ require special attention, but a
	similar approach.

\end{ex}

\begin{ex}\
	\begin{enumerate}

		\item Suppose that $G$ is a finite group, with $H$ and $K$
		      subgroups of $G$ such that $H \cap K = \{e\}$, $H \subseteq N_G(K)$ and
		      $K \subseteq N_G(H)$.  Prove that $HK \cong H \times K$.

		\item Suppose $G$ is a finite group, and suppose all of its Sylow subgroups are normal.
		      Prove that $G$ is a direct product of its Sylow subgroups.

	\end{enumerate}

\end{ex}

\begin{ex}
	Prove that every group of order $3\cdot11\cdot41$ is cyclic.

\end{ex}

\begin{ex}
	Prove that every group of order $312 = 2^3\cdot3\cdot13$ has a normal Sylow subgroup.
\end{ex}
\begin{proof}
	Consider that $N_{13} \equiv 1 \bmod 13$, and $N_{13}$ divides $2^3\cdot 3 = 24$.
  The only possible values for $N_{13}$ less than 24 are 1 and 14, but 14 doesn't divide 24, so $N_{13} = 1$, and thus the one Sylow 13-subgroup is normal.
\end{proof}

\begin{ex}
	Suppose $G$ is a group of order 105.  Prove that $G$ has a normal subgroup
	of order 35, and prove that if it also has a normal Sylow 3-subgroup then $G$ is cyclic.
\end{ex}
\begin{proof}
  Note that $105 = 3 \cdot 5 \cdot 7$.
  Then $N_7$ divides 15, so since $N_7 \equiv 1 \bmod 7$, either $N_7 = 1$ or $N_7 = 15$.
  If $N_7 = 15$, then each of the 15 7-subgroups must intersect only at the identity, so this would account for $15 \cdot 6 + 1 = 91$ elements, leaving only 14 elements to be taken up by 3-subgroups and 5-subgroups.
  We observe that by similar order considerations, $N_5$ is either 1 or 21.
  Clearly $N_5$ cannot be 21 since each such subgroup would have 4 unique elements, and $4 \cdot 21 > 14$.
  Thus $N_5 = 1$, meaning we have now accounted for 95 elements.
  Then $N_3$ can either be 1 or 7, but it cannot be 7 since we don't have $7 \cdot 2 = 14$ more elements to work with, so we must have $N_3 = 1$,
  But then there exists a normal subgroup of order 15, which must be cyclic as shown in the book, so that adds 14 more elements that we haven't yet accounted for (since they are all of order 15).
  This finally contradicts the assumption that $N_7 = 15$, so we conclude that $N_7 = 1$.

  Taking this into account, we have 7 elements accounted for to start with.
  We now have $N_5$ being either 1 or 21 again.
  If $N_5 = 21$, then that accounts for $21 \cdot 4 = 84$ more elements, bringing us to a total of 91.
  Now we consider whether $N_3$ is 1 or 7.
  If $N_3 = 7$, then we have another 14 elements to add into the mix, which brings us to exactly 105.
  I'm not sure where to go from there since $N_5$ seems like it can be 21 and thus a 5-subgroup may not be normal.
\end{proof}

\begin{ex}
	Let $G$ be a finite group, let $p$ be a prime dividing $|G|$, and let $P$ be
	a Sylow $p$-subgroup of $G$.  Prove\footnote{Note that if $P$ is a normal subgroup
		of another subgroup $N$ of $G$, then it is $p$-Sylow in $N$.} that $N_G(N_G(P)) = N_G(P)$.

\end{ex}

\begin{ex}\
	\begin{enumerate}

		\item Let $G = \GL_n(\mathbb{F}_p)$.  Let $H$ be the set of upper triangular matrices
		      with 1 on the diagonal.  Prove that $H$ is a Sylow $p$-subgroup of $G$.

		\item Determine\footnote{First, determine the possibilities for $n_p$, and then
			      show that a Sylow $p$-subgroup is not normal.} $n_p$ for $\GL_2(\mathbb{F}_p)$.

	\end{enumerate}
\end{ex}

\end{document}
